\section{Related Work}
Reliable energy benchmarking has two distinct requirements: comparable experiments and measurements that preserve the quantity of interest. MLPerf Tiny and MLPerf Power address the first by defining workloads and energy-reporting procedures, while MLPerf Inference specifies power-measurement rules \cite{banbury2021mlperftiny,tschand2024mlperfpower,mlcommonsInferencePower}. Edge studies such as DeepEdgeBench and MLPerf-based USB-accelerator benchmarking use such measurements to compare platform behaviour \cite{baller2021,libutti2020usbaccelerators}. These protocols provide the experimental context; selecting a rate for a particular integration interval remains a separate measurement decision.

External instrumentation makes that decision testable. Jetson sensor calibration against external measurements addresses systematic agreement \cite{shalavi2023jetsoncalibration}, and PowerSensor3 provides an open acquisition system with up to 20~kS/s \cite{vanderVlugt2025powersensor3}. Most directly related, Gralka et al.'s \emph{Power Overwhelming} uses synchronised oscilloscope acquisition to compare meters and resolve structure within rendered frames \cite{gralka2024poweroverwhelming}. Their distinction between aggregate agreement and temporal detail is important here: a useful energy estimate need not expose short-lived activity. We build on that distinction by quantifying rate--duration choices for energy and evaluating a separate spectral-coverage criterion on edge-AI executions.

Telemetry introduces another distinction: polling an interface faster need not produce a faster physical observation. Studies of NVIDIA sampling behaviour and software-based meters expose limitations of the underlying sensors and interfaces \cite{yang2024parttime,jay2023softwarepowermeters}. Dauner et al. directly study measurement frequency for RAPL and NVML, finding that update behaviour and acquisition overhead affect the interpretation of energy readings \cite{dauner2026frequency}. Their controlled variable is the interface-read interval. In our same-trace analysis, the physical execution is unchanged; the variable is the sampling grid used to reconstruct its energy. Direct acquisitions then test what this numerical comparison leaves out.

The contribution is thus a connection between instrument validation and benchmark practice. Rather than proposing a universally sufficient meter rate, we relate energy error and measured fluctuation coverage to a stated interval, then check execution variability and measurement boundaries. This makes explicit which objective a selected rate supports and where deployment validation is still needed.
